%% ============================================================================
\section{Introduction}
\label{ch:introduction}
\addcontentsline{lof}{section}{Introduction}
%% ============================================================================

\subsection{Background and Motivation}

Remote sensing has become a central instrument of archaeological prospection.
Where a field survey of a large region once required years of systematic ground
inspection, openly available satellite imagery now permits candidate sites to be
identified from a desk, and the resulting shortlist to direct where limited
excavation effort is spent~\cite{Casini23,Orengo20}.
The southern Mesopotamian floodplain is a canonical setting for this: it contains
thousands of mounded settlement sites, locally called \emph{tells}, formed by the
accumulated building levels and occupation debris of communities that remained in
one place over long periods~\cite{Wilkinson03}.
Their colour and shape make them visible from above, and that visibility is what
makes automated detection plausible at all.

Semantic segmentation - assigning a class label to every pixel - is well matched
to this task, because archaeologists need the outline of a site rather than merely
its presence.
Casini et al.~\cite{Casini23} fine-tuned pretrained segmentation networks on a corpus
of roughly five thousand ground-truthed sites and reported intersection-over-union
(IoU) in the range of 74--83\%, with site-level detection accuracy near 80\% once a
domain expert had adjudicated the predictions.
That work also established the framing that this thesis inherits: the model is a
component in a human--AI workflow, not a replacement for the archaeologist.

Figure~\ref{fig:prediction-examples} shows what the task looks like in practice,
using the reproduction checkpoints of Appendix~\ref{app:reproduction} on
held-out tiles. Where a mound is tonally distinct from the surrounding land, all
three architectures recover its outline closely and agree with one another;
performance falls away on small sites, and
Section~\ref{sec:repro-qualitative} characterises that dependence in full.

\begin{figure}[!ht]
\centering
\includegraphics[width=\linewidth]{prediction-examples.png}
\caption[Example predictions on held-out test tiles]{The segmentation task.
Columns: input satellite imagery; the archaeologist's annotated outline
(filled, solid boundary); and each architecture's predicted outline (filled,
solid boundary) with the annotation repeated as a dashed line for comparison.
Rows one to three are representative of the regime where all three
architectures agree; row four is a characteristic failure on a small site; row
five is a tile containing no site. Examples are selected, not randomly
sampled -- Section~\ref{sec:repro-qualitative} gives the selection basis and
the full score distribution.}
\label{fig:prediction-examples}
\end{figure}

The present work takes that workload as given and asks a different question.
Training such models on a shared HPC cluster is conventionally treated as a problem
of \emph{accelerator} throughput -- the accelerator being the specialised
processor that performs the training computation, in this work an NVIDIA GPU
throughout, and the term used where the argument does not depend on the device
being a GPU specifically.
The model is profiled, the batch size is tuned, mixed
precision is enabled, and the data path is assumed to keep up.
For this workload that assumption does not hold.
The corpus consists of thousands of individually stored image and mask files on a
shared parallel filesystem, and each training step issues an independent
\code{open}/\code{stat}/\code{read}/\code{close} sequence per sample against a
filesystem optimised for large sequential transfers rather than for many-small-file
metadata operations.
The accelerator waits.

\subsection{The problem}

Three properties of this setting, taken together, define the problem addressed here.

\paragraph{The bottleneck is real and it is not the model.}
Profiling the baseline pipeline across three architectures shows a ratio of
data-path time to compute time large enough that the loader is a plausible
constraint rather than a rounding error, and the effect is most pronounced for
the architecture with the lowest per-step compute - precisely the configuration
in which a fast model exposes a slow pipeline.

\paragraph{The bottleneck moves, but the configuration does not.}
\code{num\_workers}, \code{prefetch\_factor} and \code{pin\_memory} are fixed at
loader construction and never revisited.
On a shared filesystem, observed read latency is not a property of the dataset
alone but of whatever else is running on the cluster, and therefore varies over the
lifetime of a single job.
A statically configured loader cannot respond to a bottleneck that shifts underneath it.

\paragraph{Future accesses are known exactly, not merely predictable.}
Once the epoch permutation is drawn, the sequence of samples - and hence the
sequence of files - is fully determined.
Prefetching in this setting need not be speculative.
This is the property that the method exploits, and it is what separates an exact
staging schedule from a heuristic cache.

\subsection{Research question}

The thesis addresses one question:

\begin{quote}
Can a profiler-adaptive \code{DataLoader} that combines bottleneck classification,
tar-shard staging, and predictive prefetching improve GPU utilisation and tile
throughput for large-scale satellite-image segmentation on HPC systems?
\end{quote}

which decomposes into four sub-questions, each answered in this work by a specific
experimental contrast rather than by an aggregate result:

\begin{description}
  \item[RQ1.] When does the \code{DataLoader} become the bottleneck on a shared
    parallel filesystem?
  \item[RQ2.] Do tar shards reduce many-small-file overhead?
  \item[RQ3.] Does predictive scratch staging outperform fixed shard streaming?
  \item[RQ4.] Can prefetch depth be auto-tuned from measured timings?
\end{description}

\subsection{Contributions}

\begin{enumerate}
  \item \textbf{A characterisation of the workload as an I/O problem.}
    We describe the segmentation corpus not only by its class balance and evaluation
    metrics, as is conventional, but by its file count, size distribution and access
    pattern - the characterisation that motivates staging and that the segmentation
    literature usually leaves implicit (Section~\ref{ch:materials}).

  \item \textbf{\pads{}, a profiler-adaptive data-staging loader.}
    A drop-in \code{DataLoader} replacement that measures \Tdata{}, \Tdecode{},
    \Thtod{} and \Tgpu{} over a bounded warm-up window, classifies the dominant
    bottleneck by an inspectable rule set, selects among loose, shard and stage
    policies, and derives its prefetch depth from the measured staging and
    consumption times rather than from a user-supplied constant
    (Section~\ref{ch:methodology}).

  \item \textbf{An evaluation that isolates each mechanism.}
    \pads{} is compared against an unmodified loader, a manually tuned loader,
    fixed tar-shard streaming, and a full pre-staging upper bound, so that the
    benefit of sharding is separated from the benefit of staging, and the benefit
    of staging from the benefit of \emph{predicting} what to stage
    (Section~\ref{ch:results}).

  \item \textbf{A reproduction of the reference segmentation results, and an account
    of what it takes to obtain them.}
    We reproduce the reference configuration of Casini et al.~\cite{Casini23} and
    document the pipeline defects and environmental sensitivities encountered in
    doing so, including a case in which reduced batch size - forced by accelerator
    memory - degrades IoU substantially through its effect on batch normalisation
    statistics. That diagnosis is then confirmed directly rather than left as an
    inference: a full-protocol replica at the reference's own batch size, epoch
    budget and final-epoch weighting closes the gap to within $1$--$1.4$ points of
    the published figure, across all three architectures used throughout this
    thesis (Appendix~\ref{app:reproduction}).
\end{enumerate}

\subsection{Thesis structure}

Section~\ref{ch:literature} reviews segmentation architectures, deep learning for
archaeological prospection, and prior work on I/O for large-scale training.
Section~\ref{ch:materials} describes the dataset in both of its characterisations
and the three model architectures.
Section~\ref{ch:methodology} presents \pads{}.
Section~\ref{ch:results} reports the evaluation, and Section~\ref{ch:ablation}
isolates the contribution of individual components.
Section~\ref{ch:conclusions} concludes, and Section~\ref{ch:futurework} discusses
limitations and further work.
